% !PRINTED 127
\weberpage{127}
%<<CONT
to which the following is still to be remarked: the roots $\sqrt{\alpha+\beta\omega}$,
$\sqrt{-i(\alpha+\beta\omega)}$ are to be taken with \gesperrt{positive} real part.
That one of them should be purely imaginary is excluded by the assumption
that $\omega$ has a positive imaginary part and that $\alpha$, or $\beta$
respectively, is positive; for accordingly $\alpha+\beta\omega$ or $-i(\alpha+\beta\omega)$
cannot be real and negative; $\left(\dfrac{\beta}{\alpha}\right)$ and
$\left(\dfrac{\alpha}{\beta}\right)$ is the \gesperrt{Legendre-Jacobi} symbol
from the theory of quadratic residues, with the extension that
$\left(\dfrac{\beta}{1}\right)$ and $\left(\dfrac{0}{1}\right)=1$ shall hold. If
in the first case $\alpha$, or in the second $\beta$, is negative, then on the right
all the signs of $\alpha, \beta, \gamma, \delta$ must be reversed.
If both $\alpha$ and $\beta$ are odd, then (15), 1. as well as (15), 2.
may be applied, and both yield, as one easily proves on the basis
of the law of reciprocity for quadratic residues, the same
 result. For, if $\alpha$ and $\beta$ are odd,
$\alpha$ is taken positive, and the upper or lower sign
holds according as $\beta$ is positive or negative, we have:
\[
\left(\frac{\beta}{\alpha}\right)\left(\frac{\pm\alpha}{\pm\beta}\right)
= \pm\, e^{-\frac{\pi i}{4}(\alpha\mp1)(\beta-1)},
\]
\[
\sqrt{\mp i(\alpha+\beta\omega)} = e^{\mp\frac{\pi i}{4}}\sqrt{\alpha+\beta\omega}
\]
and the identity of (15), 1., 2. then follows from the congruences

\[
\begin{gathered}
-3\alpha\beta+\beta(\alpha+\delta)-(\beta^2-1)\alpha\gamma \equiv \alpha(\gamma-\beta)-(\alpha^2-1)\beta\delta \ (\mathrm{mod}\ 24), \\
\alpha\gamma+\alpha\beta \equiv \beta\delta \ (\mathrm{mod}\ 8),
\end{gathered}
\]
of which the second holds also for the modulus 24 whenever $\alpha$ and $\beta$
are both not divisible by 3, so that $\alpha^2\equiv\beta^2\equiv1 \ (\mathrm{mod}\ 3)$.

That the formulae (8) to (11) are satisfied by (15) is
immediately evident, and it still remains to be shown that (13)
and (14) are fulfilled. We begin with (14), where it may be
assumed that $\alpha$ is odd (and positive); for if one interchanges
$\omega$ in (14) with $-1:\omega$, then $\alpha$ and $-\beta$ are interchanged, and
these cannot both be even.

We have
\[
\sqrt{-i\omega}\ \sqrt{\alpha-\frac{\beta}{\omega}} = \sqrt{-i(-\beta+\alpha\omega)};
\]
for if we set for the moment
%CONT>>
